\section{A scalar identity that controls area and mass}
\label{sec:scalar}
\label{sec:scalar:construction}

Throughout this section the data satisfy Definition~\ref{def:prepared}.  We seek a metric
$\widehat g=e^{2t}(g+l(t)^2df\otimes df)$ with
$t>-\epsilon$, so that $\widehat g\ge e^{-2\epsilon}g$ controls
the loss of full-cut area.  The normalized height $h=\eta f$ will
separate the filling from the enclosure; weighted coercivity will
control its derivatives.  We solve for the graph height $f$ and the
modified conformal variable $Z$ in Equation~\eqref{eq:scalar:implicit}.
The scalar divergence equation for $Z$ controls the penalty's
contribution to the ADM energy.  Section~\ref{sec:elliptic} constructs
smooth solutions by exhaustion at fixed $\epsilon,N$; the later limit
in $N$ uses only the explicit estimates.  All identities below hold
in every dimension $n=k+1\geq3$.

\subsection{Filling the boundary and choosing the unknowns}

We need a divergence equation with no inner flux. A smooth filling permits
that equation to be solved on a complete manifold without boundary. The
filling is auxiliary: its energy condition will never be assumed. Its
height will later be separated from the Riemannian comparison region.


The original exterior has compact complement of its coordinate end.
Cut it at a large coordinate sphere, take an orientation-reversed copy
of the resulting compact manifold, and cap the copied spherical face by
a ball.  The resulting compact manifold has boundary a copy of the
\emph{whole} $S$.  Gluing along that boundary therefore fills all
components of $S$ simultaneously; no individual component is required
to bound.  Smooth collars perform the gluing.  Initially extend $g,K$
across $S$ on a fixed short collar.  Strict negativity of the expansions
of the signed-distance leaves persists there by continuity.  Beyond
that collar interpolate the metric to a product, and subtract a smooth
nonnegative multiple of the metric from $K$.  Start this subtraction
while the strict inequality still holds and make the multiple large
before performing the remaining interpolation.  Since the transition
is fixed and compact, its leaf mean curvatures and the other tangential
traces are bounded.  Thus the subtraction can ensure
\[
 H+\tr_{\mathrm{leaf}}K<-c_0<0
\]
through every transition.  Arrange $K=-L_0g$ on the ensuing product
necks and on the remaining cap, with $L_0$ a sufficiently large fixed
constant.  The neck length at each boundary component may be chosen
between $c(1+N)$ and $C(1+N)$, increasing $C$ when the height ramp in
Section~\ref{sec:height} is chosen.

Denote the resulting complete manifold without boundary by
$\mathcal M_N$, and continue to write $g,K$ for the extended fields.
They are unchanged on the original exterior.  All the local metric
models belong to a fixed finite family, apart from translation along
product cylinders.  Consequently their injectivity radii are bounded
below, every prescribed finite collection of local geometric bounds
is bounded independently of $N$, and the enlarged core has volume and
covering number $O(1+N)$.  No energy condition is imposed in the added
region.  Extend the end radius to a smooth positive function $\varrho$,
uniformly comparable to $1$ on that core.  Hence weighted integrals
over the core cost at most a polynomial in $N$.

In the scalar construction, $p=p(t)$ is the cutoff defined below;
the dimensional exponent is written $n-2=k-1$.
Fix $0<\epsilon<1$.  Let $\vartheta\in C^\infty(\R)$ be nonincreasing,
equal to $1$ on $(-\infty,0]$ and to $0$ on $[1,\infty)$, with values
in $[0,1]$.  For an integer $N>\max(k,2/\epsilon)$ put
\begin{equation}
 \begin{gathered}
 p(t)=N\vartheta(Nt),\qquad
 l(t)=\exp\left(\int_0^t p(s)\,ds\right),\\
 \ell=e^{-N\epsilon},\qquad \eta=\ell^{3/2},\qquad
 m_0(t)=\vartheta\bigl(N(t+\epsilon)\bigr).
 \end{gathered}
 \label{eq:scalar:parameters}
\end{equation}
In particular $l=e^{Nt}$ for $t\leq0$, $l\leq e$ everywhere, and
$\sup|p^{(j)}|\leq C_jN^{j+1}$ for each fixed $j$.

For smooth unknowns $f,Z$, define $t$ implicitly by
\begin{equation}
 Z=t+\frac{1}{2k}\log D,\qquad D=1+l(t)^2|df|_g^2.
 \label{eq:scalar:implicit}
\end{equation}
For fixed $df$ the derivative of the right side with respect to $t$
is $(k+pv)/k\geq1$, where $v=1-D^{-1}$.  Its limits as $t$ tends to
$\pm\infty$ are $\pm\infty$.  Thus this equation defines $t$ globally
and smoothly as a function of $(Z,|df|^2)$, without a lower-bound
assumption on $t$.  Use $g$ to identify vectors and covectors and set
\begin{equation}
 \begin{aligned}
 &w=lD^{-1/2}\nabla f,\qquad a_w=|w|,\qquad v=a_w^2,\qquad d=D^{-1}=1-v,\\
 &\chi=\frac{kd}{k+pv},\qquad
 A_d=I-w\otimes w,\qquad
 A_\chi=I-\frac{k+p}{k+pv}w\otimes w,\\
 &H^f=lD^{-1/2}\Hess_g f,\qquad L=K+H^f,\qquad
 F=\tr_{A_\chi}L,\qquad h=\eta f,\\
 &U=l\sqrt D,\qquad u=e^{(k-1)t}U,\qquad
 V=uA_\chi\bigl(k\nabla Z+K(w,\cdot)\bigr),\\
 &\bar g=g+l^2df\otimes df,\qquad \widehat g=e^{2t}\bar g,
 \qquad s_p=p+\frac{k-1}{2}.
 \end{aligned}
 \label{eq:scalar:variables}
\end{equation}
The notation $\tr_A B$ means $A^{ij}B_{ij}$, and
$|\alpha|_A^2=A^{ij}\alpha_i\alpha_j$.  The tensors $A_d,A_\chi$
have transverse eigenvalue $1$ and axial eigenvalues $d,\chi$,
respectively.  More explicitly,
\begin{equation}
 0<\frac{k}{k+N}d\leq\chi\leq d\leq1,
 \qquad
 k\,dZ=(k+pv)\,dt+H^f(w,\cdot).
 \label{eq:scalar:differential}
\end{equation}
The second equality follows by differentiating
Equation~\eqref{eq:scalar:implicit}:
$\tfrac12d\log D=pv\,dt+H^f(w,\cdot)$.
Every tensor in Equation~\eqref{eq:scalar:variables} is smooth at
$df=0$.  The axis notation used below is only a way of displaying
these tensors, not a coefficient definition at gradient zeros.

In this and the next two sections $\mathcal P_N$ denotes a polynomial
in $1+N$, whose coefficients may depend on the fixed prepared data,
$n$, and $\epsilon$.  Increasing such a polynomial finitely many times
does not change this convention.  An estimate with an arbitrary
constant depending on fixed $N$ will be explicitly identified as such;
it cannot be substituted into a later $\mathcal P_N$ estimate.

\subsection{The scalar identity}

The model for this calculation is the Schoen--Yau scalar-curvature
identity for Jang graphs \cite{SchoenYau1981}, and more specifically
its warped generalization in \cite[Identity~9, p.~580]{BrayKhuri2011}.
Our conformal variable, gradient-aligned trace, and redistributed
divergence are different. We derive their exact identity before
imposing the coupled equations; its algebra and their global
solvability are separate steps.

For symmetric tensors define the bilinear form and trace reversal
\[
 \mathsf B(B,C)=\langle B,C\rangle_g-(\tr_gB)(\tr_gC),
 \qquad P_B=B-(\tr_gB)g.
\]
At a nonzero gradient let $e=\nabla f/|df|$, and decompose $L$ and $dt$
relative to $e^\perp\oplus\R e$ as
\[
 L=\begin{pmatrix}T&M\\M^\intercal&j\end{pmatrix},\qquad dt=(y,x),
 \qquad T^0=T-\frac{\tr T}{k}I_{e^\perp}.
\]
Here $\tr T=F-\chi j$.  At $df=0$ one may choose any unit axis.

\begin{proposition}[Scalar identity]\label{prop:scalar:identity}
For every smooth $f,Z$ and every smooth $g,K$ the following identity
holds, without a field equation or energy hypothesis:
\begin{equation}
 \frac12e^{2t}R_{\widehat g}
 =\mu+J(w)+\mathcal T+w(F)-F\tau-u^{-1}\divg_g V,
 \label{eq:scalar:identity}
\end{equation}
where $2\mu=R_g+\tau^2-|K|^2$, $J=\divg_gP_K$, and
\begin{equation}
 \begin{aligned}
 \mathcal T={}&\frac12|T^0|^2+|M+s_pa_wy|^2
       +\left(ks_p-\frac{(k-1)^2}{4}v\right)|y|^2+ks_pd x^2\\
 &-\frac{k-1}{2k}(F-\chi j)^2+\chi j^2-\chi Fj+F^2
       +2s_p\chi j a_w x .
 \end{aligned}
 \label{eq:scalar:quadratic}
\end{equation}
The quantity $\mathcal T$ is smooth, including where $df=0$.
\end{proposition}

\begin{proof}
We derive the curvature formula before completing any squares.  On
$\R_z\times\mathcal M_N$ consider the stationary Lorentz metric
\[
 \mathbf G=-l^2(dz-df)^2+\bar g
          =-l^2dz^2+2l^2dz\,df+g.
\]
The constant-$z$ slices have lapse $U$, shift $Uw=l^2\nabla f$,
future unit normal $\mathbf n=U^{-1}\partial_z-w$, and second
fundamental form
\begin{equation}
 b=-U^{-1}\sym\nabla(Uw)
   =-H^f-p(dt\otimes w+w\otimes dt).
 \label{eq:scalar:stationary-second-form}
\end{equation}
This sign follows directly from
$\partial_zg=2Ub+\mathcal L_{Uw}g=0$ and the prescribed future-normal
convention.  Write $B_0=\tr b$.  Contracted Gauss and the normal lapse
equation give
\[
 \mathbf R+2\mathbf\Ric(\mathbf n,\mathbf n)=R_g-\mathsf B(b,b),
 \qquad
 \mathbf\Ric(\mathbf n,\mathbf n)
   =-\mathbf n(B_0)-|b|^2+U^{-1}\Delta_gU.
\]
Since $B_0$ is stationary, $\mathbf n(B_0)=-w(B_0)$, and taking the
trace in Equation~\eqref{eq:scalar:stationary-second-form} gives
$\divg(Uw)=-UB_0$.  Substitution proves
\begin{equation}
 \mathbf R=R_g+\mathsf B(b,b)
  -2U^{-1}\divg\bigl(\nabla U+B_0Uw\bigr).
 \label{eq:scalar:stationary-curvature}
\end{equation}

In coordinates $z'=z-f$, the same metric is the static warped product
$-l^2(dz')^2+\bar g$.  Its only mixed connection terms are
$\boldsymbol\Gamma^{z'}_{iz'}=\partial_i\log l$ and
$\boldsymbol\Gamma^i_{z'z'}=l\bar g^{ij}\partial_jl$; contracting
their curvature contributions yields
$\mathbf R=R_{\bar g}-2l^{-1}\Delta_{\bar g}l$.
Also $\bar g^{-1}=A_d$, $dV_{\bar g}=\sqrt D\,dV_g$, and therefore
\[
 l^{-1}\Delta_{\bar g}l
   =U^{-1}\divg\bigl((U/l)A_d\nabla l\bigr).
\]
Differentiating $U=l\sqrt D$ and using
Equation~\eqref{eq:scalar:differential} proves the vector identity
\[
 \nabla U-(U/l)A_d\nabla l=-Ub(w,\cdot).
\]
Consequently Equation~\eqref{eq:scalar:stationary-curvature} becomes
\begin{equation}
 \frac12R_{\bar g}
   =\frac12R_g+\frac12\mathsf B(b,b)
          +U^{-1}\divg(UP_bw).
 \label{eq:scalar:bar-first}
\end{equation}

Set $Q=K-b=L+p(dt\otimes w+w\otimes dt)$.  For any symmetric $C$,
the product rule and Equation~\eqref{eq:scalar:stationary-second-form}
give the useful exact contraction
\[
 U^{-1}\divg(UP_Cw)=(\divg P_C)(w)-\mathsf B(C,b).
\]
Expanding $\mathsf B(K-b,K-b)$ and applying this formula to $C=K$
transforms Equation~\eqref{eq:scalar:bar-first} into
\begin{equation}
 \frac12R_{\bar g}=\mu+J(w)+\frac12\mathsf B(Q,Q)
                         -U^{-1}\divg(UP_Qw).
 \label{eq:scalar:bar-identity}
\end{equation}
For example, the terms involving $K$ on the right cancel to
$R_g/2-\mathsf B(b,b)/2+(\divg P_b)(w)$, as they must.

The conformal scalar law in dimension $k+1$ is
\[
 \frac12e^{2t}R_{\widehat g}
   =\frac12R_{\bar g}-k\Delta_{\bar g}t
                         -\frac{k(k-1)}2|dt|_{A_d}^2.
\]
The volume formula gives
\[
 \Delta_{\bar g}t
  =U^{-1}\divg(UA_d\nabla t)-p|dt|_{A_d}^2.
\]
Changing the divergence weight from $U$ to
$u=e^{(k-1)t}U$ now yields
\begin{equation}
 \begin{aligned}
 \frac12e^{2t}R_{\widehat g}
  ={}&\mu+J(w)+\frac12\mathsf B(Q,Q)
     +(k-1)P_Q(w,dt)+ks_p|dt|_{A_d}^2\\
   &-u^{-1}\divg\bigl(u(P_Qw+kA_d\nabla t)\bigr).
 \end{aligned}
 \label{eq:scalar:invariant-preidentity}
\end{equation}
In particular the last change of weight adds
$k(k-1)|dt|_{A_d}^2$ to the preceding coefficient
$k[p-(k-1)/2]$.

There are two remaining first-order algebraic identities:
\begin{equation}
 \begin{aligned}
 u(P_Qw+kA_d\nabla t)&=V-uFw,\\
 u^{-1}\divg(uw)&=\tr_gL-\tau+2s_pw(t)
                  =F+(1-\chi)j-\tau+2s_pa_wx.
 \end{aligned}
 \label{eq:scalar:flux-relations}
\end{equation}
For the first, $P_Qw+kA_ddt$ has transverse and axial components
\[
 a_wM+(k+pv)y,\qquad -a_w\tr T+kd x,
\]
whereas $V/u$ has components
$a_wM+(k+pv)y$ and $\chi a_wj+kd x$ by
Equation~\eqref{eq:scalar:differential}.  Their difference is $Fw$.
For the second, use $\divg(Uw)=-U\tr b$ and change the weight.
Thus the smooth invariant definition of the claimed quadratic form is
\begin{equation}
 \mathcal T=\frac12\mathsf B(Q,Q)+(k-1)P_Q(w,dt)
        +ks_p|dt|_{A_d}^2+F\bigl(\tr_gL+2s_pw(t)\bigr).
 \label{eq:scalar:invariant-quadratic}
\end{equation}
This proves smoothness and Equation~\eqref{eq:scalar:identity}.
To obtain its displayed block form, the blocks of $Q$ are
$T$, $M+pa_wy$, and $j+2pa_wx$.  The transverse terms are
\[
 |M+pa_wy|^2+(k-1)a_w(M+pa_wy)\cdot y
 =|M+s_pa_wy|^2-\frac{(k-1)^2}{4}v|y|^2.
\]
The terms containing $F a_wx$ cancel, while their remaining axial
cross term is $2s_p\chi ja_wx$.  Finally
$|T|^2=|T^0|^2+(F-\chi j)^2/k$ gives
Equation~\eqref{eq:scalar:quadratic}.

For completeness this derivation retains derivatives of $p$.
In the flux of Equation~\eqref{eq:scalar:invariant-preidentity} the
part explicitly containing $p$ is
$p\{v\nabla t-w(t)w\}$; its divergence contributes
$p'v|y|^2$.  Thus the scalar side contains $-p'v|y|^2$.
In Equation~\eqref{eq:scalar:identity} the redistribution of this term
can also be checked directly.  Holding $t,dt,df,H^f,K$ fixed when taking
the partial derivative in $p$, one has
$\chi_p=-\chi v/(k+pv)$.  The terms containing $p'$ in $w(F)$ and
$u^{-1}\divg V$ are respectively
\[
 p'\chi_p a_wxj,
 \qquad p'\bigl(v|y|^2+\chi_p a_wxj\bigr).
\]
Their difference is exactly $-p'v|y|^2$.  No derivative of the cutoff
has been discarded.
\end{proof}

\subsection{Coercivity with polynomial constants}

The identity has isolated the divergence, but it is useful only if the
remaining quadratic form controls the graph and conformal derivatives.
The next proposition records exactly which axial weights survive. Its
constants are polynomial in the large parameter; this is what will allow
the exponentially small penalty flux to tend to zero.


\begin{proposition}[Weighted coercivity]\label{prop:scalar:coercivity}
For $0\leq p\leq N$ and $0\leq v<1$ the quadratic form in
Equation~\eqref{eq:scalar:quadratic} is nonnegative, and
\begin{equation}
 |T|^2+|M|^2+d j^2+F^2+|dt|_{A_d}^2
                 \leq C_k(1+N)^2\mathcal T.
 \label{eq:scalar:coercivity}
\end{equation}
At $dt=0$ one has, with constants depending only on $k$,
\begin{equation}
 c_k\bigl(|T^0|^2+|M|^2+F^2+\chi j^2\bigr)
 \leq\mathcal T\leq
 C_k\bigl(|T^0|^2+|M|^2+F^2+\chi j^2\bigr).
 \label{eq:scalar:stationary-coercivity}
\end{equation}
\end{proposition}

\begin{proof}
The entire axial part of Equation~\eqref{eq:scalar:quadratic} is
exactly
\begin{equation}
 \begin{aligned}
 &ks_pd\left(x+\frac{a_wj}{k+pv}\right)^2
 +\frac{k+1}{2k}\left(F-\frac{\chi j}{k+1}\right)^2
 +\frac{kd B_k(v)}{(k+pv)^2}j^2,\\
 &B_k(v)=d\frac{k(k+2)}{2(k+1)}+v\frac{k+1}{2}
        =\frac{k(k+2)+v}{2(k+1)}.
 \end{aligned}
 \label{eq:scalar:completed-axial}
\end{equation}
Indeed completing first the $x$ square and then the $F$ square leaves
\[
 \chi-\frac{k}{2(k+1)}\chi^2
             -s_p\frac{\chi^2v}{kd}
       =\frac{kd B_k(v)}{(k+pv)^2}.
\]
The last coefficient is at least $c_kd/(1+N)^2$, and
\[
 ks_p-\frac{(k-1)^2}{4}v\geq\frac{k^2-1}{4}>0.
\]
These inequalities control $|T^0|^2,|y|^2,dj^2$ and the two completed
squares.  Recovering $M$ from $M+s_pa_wy$ costs at most
$C_k(1+N)^2$.  Recovering $\sqrt d\,x$ costs the already controlled
$\sqrt d\,j$, since $a_w/(k+pv)\leq1/k$.  Recovering $F$ costs no
more because $\chi^2\leq d$.  The relation $\tr T=F-\chi j$ then
controls $|T|^2$.  These observations give
Equation~\eqref{eq:scalar:coercivity} with the stated single quadratic
polynomial; they do not remove the factors $d$.

When $dt=0$ the axial expression instead completes as
\[
 \frac{k+1}{2k}\left(F-\frac{\chi j}{k+1}\right)^2
       +\chi\left(1-\frac{k\chi}{2(k+1)}\right)j^2.
\]
Its second parenthesis is bounded above and below by positive
constants, and $\chi^2\leq\chi$.  This proves
Equation~\eqref{eq:scalar:stationary-coercivity} independently of $N$.
\end{proof}

\subsection{The penalized system and the lower barrier}

We now choose the equations using the signs in
Equation~\eqref{eq:scalar:identity}.  Imposing $F=h=\eta f$ turns
$w(F)$ into the nonnegative term $\eta a_w|df|$.  We can assign a
small fixed fraction of this term and of $\mathcal T$ to the divergence
source while retaining the rest in the scalar-curvature formula.
To enforce the lower bound for $t$, the source will also contain a
negative term near $t=-\epsilon$.  Its sign improves scalar curvature,
but its integral can increase ADM energy.  We therefore need both a
pointwise floor argument and a separate estimate for the weighted
integral of that penalty.

Choose
\begin{equation}
 \frac n4<b_1<\min\left(n-2,\frac n2\right),\qquad
 \rho_0=\varrho^{-n-\beta},\quad \rho_1=\varrho^{-2b_1},\quad
 \rho=c_*\rho_0>0.
 \label{eq:scalar:weights}
\end{equation}
This interval is nonempty for every $n\geq3$.  Its three restrictions
serve different parts of the argument: $b_1<n-2$ permits the radial
height comparison in Section~\ref{sec:elliptic}; $2b_1<n$ lets
$v\rho_1$ dominate the end-curvature error in the floor calculation;
and $4b_1>n$ makes $\rho_1^2$ integrable in the flux estimate.
Take $c_*$ small enough that $\mu-|J|\geq2\rho$ on the original
exterior.  The system to be solved is
\begin{equation}
 \begin{gathered}
 F=h=\eta f,\qquad \divg V=u\Xi,\\
 \Xi=\delta_0\bigl(\mathcal T+\eta a_w|df|\bigr)+\rho
                  -m_0 C_N(\rho_0+v\rho_1).
 \end{gathered}
 \label{eq:scalar:system}
\end{equation}
Here $\delta_0>0$ is a sufficiently small fixed number and $C_N$ is
a sufficiently large polynomial, specified by the next lemma.
Substitution in the scalar identity retains
$(1-\delta_0)(\mathcal T+\eta a_w|df|)$; the positive source $\rho$
spends only part of the strict energy margin on the original exterior.
The remaining term $-h\tau$ will be controlled on a small height band.
The negative penalty is chosen to contradict the flux lower bound at
a minimum with $t=-\epsilon$.

For an above-floor solution its support lies in
$-\epsilon<t<-\epsilon+N^{-1}$, where $\ell\le l\le e\ell$.
Proposition~\ref{prop:scalar:flux} will use this small weight to bound
the possible energy increase by a polynomial times
$\ell+\ell^2/\eta$.  The height-derivative estimate below instead
uses $\eta/\ell$.  The choice $\eta=\ell^{3/2}$ makes both ratios
$\ell^2/\eta$ and $\eta/\ell$ equal to $\ell^{1/2}$.  Both ratios
tend to zero faster than any polynomial in $N$ at fixed $\epsilon$.

On a large truncation the outer data are $f=Z=0$; there is no inner
boundary.  For the existence argument we use two homotopies.  In the
first replace $K$ by $aK$, $0\leq a\leq1$, everywhere in these
definitions, keeping the weights and constants fixed.  In the second
keep $K=0$ and multiply the whole expression $\Xi$ by $s\in[0,1]$.

For a symmetric tensor $B$ write
\[
 \mathcal S(B)=\frac12\left\{(\tr_{A_d}B)^2
                                      -|B|_{A_d\otimes A_d}^2\right\}.
\]

\begin{lemma}[Pointwise exclusion of the floor]\label{lem:scalar:floor}
There is $c_1>0$ depending on $k$ alone such that, at any point with
$dt=0$,
\begin{equation}
 \mathcal T-\mathcal S(H^f)
       \geq c_1\mathcal T-\mathcal P_N(vF^2+|K|^2).
 \label{eq:scalar:floor-algebra}
\end{equation}
Suppose a smooth solution on a truncation with outer $f=Z=0$
satisfies the trace equation and has $|h|\leq C$ on
the core and $|h|\leq C\varrho^{-b_1}$ on the end, with constants
independent of the homotopy parameter and $N$.  One may choose
$0<\delta_0<\min(c_1,1)/2$ and a polynomial $C_N$ such that a solution
of Equation~\eqref{eq:scalar:system} for which $\min t=-\epsilon$
cannot attain this minimum of $t$ at an interior point.  This holds in
both stages of the homotopy.
\end{lemma}

\begin{proof}
At $dt=0$, direct block contraction gives
\[
 \mathcal S(L)=\frac{k-1}{2k}(F-\chi j)^2
      -\frac12|T^0|^2+dj(F-\chi j)-d|M|^2.
\]
Subtracting from Equation~\eqref{eq:scalar:quadratic} gives the full
difference, including its mixed term:
\begin{equation}
 \begin{aligned}
 \mathcal T-\mathcal S(L)={}&|T^0|^2+(1+d)|M|^2+\frac{F^2}{k}\\
  &+\left(\frac{k-2}{k}\chi-d\right)Fj
        +\chi\left(1+d-\frac{k-1}{k}\chi\right)j^2.
 \end{aligned}
 \label{eq:scalar:floor-difference}
\end{equation}
If $(1+p)v\leq\zeta_k$ with $\zeta_k>0$ sufficiently small, then
$|d-1|+|\chi-1|\leq C_k\zeta_k$.  At $d=\chi=1$ the expression is
\[
 |T^0|^2+2|M|^2+\frac{(F-j)^2}{k}+j^2=|L|^2.
\]
Perturbation of this finite-dimensional positive form, and
Equation~\eqref{eq:scalar:stationary-coercivity}, therefore bound
Equation~\eqref{eq:scalar:floor-difference} below by a fixed positive
multiple of $\mathcal T$.

In the other case $1/v\leq(1+p)/\zeta_k$.  The last coefficient in
Equation~\eqref{eq:scalar:floor-difference} is at least $\chi$,
since $\chi\leq d$.  The absolute value of its mixed coefficient is
at most $d$.  Hence Young's inequality bounds the difference below by
\[
 |T^0|^2+|M|^2+\tfrac12\chi j^2
                  +\left(\frac1k-\frac{d^2}{2\chi}\right)F^2.
\]
Here
\[
 \frac{d^2}{\chi}=\frac{d(k+pv)}k\leq1+\frac Nk.
\]
It follows, again using
Equation~\eqref{eq:scalar:stationary-coercivity}, that the desired
fixed multiple of $\mathcal T$ is obtained after subtracting at most
$C_k(1+N)^2vF^2/\zeta_k$.

To replace $L$ by $H^f=L-K$, polarize $\mathcal S$ in the
$A_d\otimes A_d$ norm.  The preceding stationary coercivity gives
\[
 |L|_{A_d\otimes A_d}^2
   =|T|^2+2d|M|^2+d^2j^2\leq C_k(1+N)\mathcal T.
\]
Moreover $|K|_{A_d\otimes A_d}\leq|K|$.  Therefore
\[
 |\mathcal S(L)-\mathcal S(L-K)|
       \leq \zeta\mathcal T+C_{k,\zeta}(1+N)|K|^2
\]
for every fixed $\zeta>0$.  Taking $\zeta$ below half the positive
constant already obtained proves
Equation~\eqref{eq:scalar:floor-algebra}.

At an interior minimum of $t$, $dt=0$ and $\Hess_g t\geq0$.
To estimate curvature from above, view $\bar g$ as the metric on the
graph of $f$ in the Riemannian warped product $g+l^2dz^2$.
At that point $dl=0$ and $l^{-1}\Hess l=p\Hess t$.  Its ambient
scalar curvature is $R_g-2p\Delta t$, its horizontal Ricci tensor
is $\Ric_g-p\Hess t$, and its unit vertical Ricci component is
$-p\Delta t$.  The graph unit normal is
$D^{-1/2}\partial_{z,\mathrm{unit}}-w$ and its second form with this
orientation is $-H^f$, which has the same quadratic Gauss term as $H^f$.
The Gauss equation followed by the conformal law consequently gives
the exact minimum-point formula
\begin{equation}
 \begin{aligned}
 \frac12e^{2t}R_{\widehat g}
 ={}&\frac12R_g-v\Ric_g(e,e)+\mathcal S(H^f)\\
   &-pv\tr_{e^\perp}\Hess t-k\tr_{A_d}\Hess t\\
 \leq{}&\frac12R_g-v\Ric_g(e,e)+\mathcal S(H^f).
 \end{aligned}
 \label{eq:scalar:minimum-curvature}
\end{equation}
When $v=0$, terms using $e$ vanish.  In particular all $p'(dt)^2$
terms vanish at this point.

Combining Equations~\eqref{eq:scalar:identity},
\eqref{eq:scalar:floor-algebra}, and
\eqref{eq:scalar:minimum-curvature}, and using $F=\eta f$, yields
\begin{equation}
 u^{-1}\divg V\geq c_1\mathcal T+\eta a_w|df|
                              -\mathcal P_N(\rho_0+v\rho_1).
 \label{eq:scalar:minimum-flux}
\end{equation}
Here are the weights in this estimate.  The curvature-independent
constraint contribution is
$\mu-R_g/2=(\tau^2-|K|^2)/2$.  This, $J(w)$, $F\tau$, and the
$|K|^2$ error are uniformly bounded on the enlarged core and vanish
on the far end, so they cost $C\rho_0$ pointwise.  On the end,
$v|\Ric_g|\leq Cv\varrho^{-n}\leq Cv\rho_1$ because $2b_1<n$;
the height comparison gives $vF^2\leq Cv\rho_1$.  The same statements
hold uniformly for $aK$, $0\leq a\leq1$.  No use of the dominant
energy condition in the filling has entered this calculation.

At $t=-\epsilon$ one has $m_0=1$.  Subtracting the right side of
Equation~\eqref{eq:scalar:system} from the lower bound in
Equation~\eqref{eq:scalar:minimum-flux} gives
\[
 (c_1-\delta_0)\mathcal T+(1-\delta_0)\eta a_w|df|
        +(C_N-\mathcal P_N)(\rho_0+v\rho_1)-\rho.
\]
Choose $C_N>\mathcal P_N+c_*+1$.  The expression is strictly
positive, contradicting the equation.  In the second homotopy stage
the trace comparison gives $f=0$: a positive maximum or negative
minimum contradicts $\tr_{A_\chi}H^f=\eta f$.
Thus $t=Z$, and at a floor minimum the left side of the divergence
equation divided by $u$ is $k\Delta t\geq0$, whereas its right side
is $s(\rho-C_N\rho_0)<0$ for $s>0$.  At $s=0$, the equation
$\divg(u\nabla Z)=0$ with zero outer data gives $Z=0$ by the maximum
principle.  This finishes the pointwise exclusion.
\end{proof}

The lemma excludes the boundary $\min t=-\epsilon$ of the moving
degree domain.  It does not purport to classify solutions that lie
below that domain.  The uniform height comparisons and the separate
outer-face exclusion used in its hypotheses are proved in
Section~\ref{sec:elliptic}.

\subsection{Scalar curvature on solutions and height derivatives}

On the original exterior, substituting
Equation~\eqref{eq:scalar:system} into the identity gives
\begin{equation}
 \begin{aligned}
 \frac12e^{2t}R_{\widehat g}
  ={}&\mu+J(w)-\rho-h\tau
       +(1-\delta_0)\mathcal T
       +(1-\delta_0)\eta a_w|df|\\
    &+m_0C_N(\rho_0+v\rho_1).
 \end{aligned}
 \label{eq:scalar:solution-curvature}
\end{equation}
In particular a small fixed height band has a strict scalar margin:
choose $b_3>0$ such that $b_3|\tau|\leq\rho/2$ on the support of
$\tau$.  Then $\mu-|J|-\rho-|h\tau|\geq\rho/2$ on
$|h|\leq b_3$.  This is possible since $\tau$ is compactly supported
and $\mu-|J|\geq2\rho$ there.

For later use the exact height identities are
\begin{equation}
 \begin{aligned}
 |dh|_{A_d}&=\frac{\eta|df|}{\sqrt D}\leq\frac\eta l,\\
 e^{2t}\Delta_{\widehat g}h
   &=\frac{\eta}{l\sqrt D}
       \left\{\divg w+(k-1-p)w(t)\right\},\\
 \divg w&=\tr_{A_d}H^f+pd\,w(t).
 \end{aligned}
 \label{eq:scalar:height-derivatives}
\end{equation}
The middle formula follows from
$\sqrt D A_d\nabla h=(\eta/l)w$ and the conformal Laplacian law;
the last follows by differentiating $w=lD^{-1/2}\nabla f$.
For an above-floor solution, $l\geq\ell$, and
Proposition~\ref{prop:scalar:coercivity} therefore implies
\begin{equation}
 |dh|_{A_d}\leq\ell^{1/2},\qquad
 |e^{2t}\Delta_{\widehat g}h|
                  \leq\mathcal P_N\ell^{1/2}(1+\sqrt{\mathcal T}).
 \label{eq:scalar:height-estimates}
\end{equation}
Indeed $\tr_{A_d}H^f=\tr T+dj-\tr_{A_d}K$ is bounded by
$\mathcal P_N(1+\sqrt{\mathcal T})$.  The other potentially large
term is $w(t)=a_wx$.  The prefactor $D^{-1/2}=\sqrt d$ converts it
to $a_w\sqrt d\,x$, precisely the weighted quantity controlled by
Equation~\eqref{eq:scalar:coercivity}.  Thus no inverse power of $d$
is lost.  Section~\ref{sec:height} uses these estimates to raise the
metric inside a compact height band.

\subsection{The penalty flux and its ADM normalization}

\begin{proposition}[Flux and energy cost]\label{prop:scalar:flux}
Fix $\epsilon$ and $N$ as above.  Suppose a smooth global solution
of Equation~\eqref{eq:scalar:system} satisfies $t>-\epsilon$,
$Z\leq\mathcal P_N$, and on the coordinate end
\[
 f=O_j(e^{-c_Nr}),\qquad t=O_j(r^{2-n}),\qquad
 \Delta_gt=O(r^{-n-\beta})
\]
for all required derivative orders, with constants allowed to depend
arbitrarily on this fixed $N$.  These hypotheses will be supplied by
Theorem~\ref{thm:elliptic:filled}.  Then
\begin{equation}
 \begin{gathered}
 \mathfrak F_N:=\lim_{r\to\infty}
       \int_{S_r}g(V,\nu_g)\,dA_g
       =\int_{\mathcal M_N}u\Xi\,dV_g,\\
 \widehat E-E=-\frac{\mathfrak F_N}{k\omega},\qquad
 \mathfrak F_N\geq-\mathcal P_N(\ell+\ell^{1/2}),\\
 \widehat E\leq E+\mathcal P_N(\ell+\ell^{1/2}).
 \end{gathered}
 \label{eq:scalar:flux-energy}
\end{equation}
The polynomial in the error uses only the fixed geometry, the weights,
and the explicit parameters, not the arbitrary fixed-$N$ derivative
constants in the stated end estimates.
\end{proposition}

\begin{proof}
For fixed $N$, the end hypotheses and
Equation~\eqref{eq:scalar:differential} give
\[
 D-1=O_j(e^{-c_N'r}),\quad Z-t=O_j(e^{-c_N'r}),\quad
 V=k\nabla t+O_N(r^{3-2n})+O(e^{-c_N'r}).
\]
The notation $O_N$ here permits arbitrary fixed-$N$ constants.
Since $N\epsilon>2$ and $t\to0$, $m_0=0$ sufficiently far out.
There $\mathcal T=ks_p|dt|^2+O(e^{-c_N'r})$, so every term in
$u\Xi$ is integrable: the gradient term is $O_N(r^{2-2n})$, and
$\rho=O(r^{-n-\beta})$.  Smoothness gives integrability on the compact
filled core.  The divergence theorem on exhausting domains has no
inner boundary and proves the first equality and existence of the
flux limit.  It also proves existence of $\lim\int_{S_r}\partial_{\nu_g}t$
directly from $\Delta_gt\in L^1$.

The graph part $l^2df^2$ contributes zero to the ADM flux by its
exponential decay.  The conformal perturbation is
\[
 (e^{2t}-1)g=2t\delta+O_1(r^{4-2n}).
\]
The linear ADM integrand of $2t\delta$ is $-2k\,\partial_i t$;
the remainder has integrated flux $O(r^{2-n})\to0$.  Euclidean and
$g$-normal fluxes differ by the same vanishing order.  Thus
\[
 \widehat E-E=-\omega^{-1}\lim_{r\to\infty}
                      \int_{S_r}\partial_{\nu_g}t\,dA_g
                =-\frac{\mathfrak F_N}{k\omega}.
\]

It remains to bound the negative part of the integral independently
of those end constants.  It is supported in the penalty band
$-\epsilon<t<-\epsilon+1/N$, where $t<0$ and
$\ell\leq l\leq e\ell$.  Put $\sigma=|df|$.  Directly from the
definitions, with constants independent of $N$ on this band,
\begin{equation}
 u\leq C(\ell+\ell^2\sigma),\qquad
 uv\leq C\ell^2\sigma,\qquad
 ua_w\sigma=e^{(k-1)t}l^2\sigma^2\geq c\ell^2\sigma^2.
 \label{eq:scalar:penalty-weights}
\end{equation}
For the second estimate use
$l\sqrt D\,v=l^3\sigma^2/\sqrt{1+l^2\sigma^2}\leq l^2\sigma$.
The nonnegative terms $u\delta_0\mathcal T$ and $u\rho$ may be
dropped, while the favorable gradient term absorbs the linear
$\sigma$ terms.  More explicitly Young's inequality gives
\[
 C C_N\ell^2\sigma(\rho_0+\rho_1)
 \leq\tfrac12\delta_0c\eta\ell^2\sigma^2
       +C' C_N^2\frac{\ell^2}{\eta}(\rho_0^2+\rho_1^2).
\]
Consequently
\[
 \mathfrak F_N\geq
  -C C_N\ell\int\rho_0
  -C' C_N^2\frac{\ell^2}{\eta}\int(\rho_0^2+\rho_1^2).
\]
All integrals use $dV_g$ on $\mathcal M_N$.  They are bounded by
$C(1+N)$: on the end $\rho_0$ is integrable and $4b_1>n$ makes
$\rho_1^2$ integrable, while on the core the weights are uniformly
bounded and its volume is $O(1+N)$.  Since
$\ell^2/\eta=\ell^{1/2}$, this proves the polynomial flux estimate.
For fixed $\epsilon>0$ the resulting error tends to zero as
$N\to\infty$.
\end{proof}

Finally, every above-floor solution satisfies the pointwise comparison
\begin{equation}
 \widehat g\geq e^{-2\epsilon}g.
 \label{eq:scalar:metric-domination}
\end{equation}
This comparison and the flux estimate are applied to each smooth
solution at fixed $N$.  They require no limiting elliptic solution as
$N\to\infty$ and no regularity assertion for a minimizing boundary.
